% Part: counterfactuals
% Chapter: introduction
% Section: strict-conditional

\documentclass[../../../include/open-logic-section]{subfiles}

\begin{document}

\olfileid{cnt}{int}{str}

\olsection{কঠোর শর্তবচন}

লুইস \emph{কঠোর শর্তবচন}~$\strictif$ প্রবর্তন করেন এবং
যুক্তি দেন যে বস্তুগত শর্তবচন নয়, এটিই নিহিতকরণের
অনুরূপ। সত্যতা-সংক্রান্ত মোডাল যুক্তিবিদ্যায়
$!A \strictif !B$-কে $\Box(!A \lif !B)$ হিসেবে
সংজ্ঞায়িত করা যায়। অতএব কোনো জগতে কঠোর শর্তবচনটি
সত্য তখন এবং কেবল তখনই, যখন সংশ্লিষ্ট বস্তুগত শর্তবচনটি আবশ্যিকভাবে
সত্য।

বস্তুগত শর্তবচনের আপাত-বৈপরীত্যের নিরিখে কঠোর শর্তবচন
কেমন কাজ করে? মিথ্যা পূর্বপক্ষযুক্ত কিংবা সত্য অনুপক্ষযুক্ত
কঠোর শর্তবচন সত্যও হতে পারে, মিথ্যাও হতে পারে। উপরন্তু,
$(!A \strictif !B) \lor (!B \strictif !A)$ বৈধ নয়।
$!A \strictif !B$, $\lnot !A \lor !B$-এর সমতুল্যও
নয়; তাই এটি সত্যমান-অপেক্ষকধর্মী নয়।

% BN-SRC-420: অননুগমনে নঞর্থকৃত কঠোর শর্তবচন চাই; বস্তুগত শর্তবচন নয়।
% BN-SRC-421: প্রদর্শিত অনুগমন ও পরের আবশ্যিকতার দাবিগুলি S5 ব্যবস্থায়।
S5 মোডাল ব্যবস্থায় আমরা পাই:
\begin{align}
  !A \strictif !B & \Entails \lnot !A \lor !B
  \text{ কিন্তু:}\\
  \lnot !A \lor !B & \Entails/ !A \strictif !B\\
  !B & \Entails/ !A \strictif !B\\
  \lnot !A & \Entails/ !A \strictif !B\\
  \lnot(!A \strictif !B) & \Entails/ !A \land \lnot !B
  \text{ কিন্তু:}\\
  !A \land \lnot !B & \Entails \lnot(!A \strictif !B)
  \intertext{তবু কঠোর শর্তবচনের জন্য মোডাস পোনেন্স চলে:}
  !A, !A \strictif !B & \Entails !B
  \intertext{কঠোর শর্তবচনে দুটি অনুপক্ষ একত্র করা যায়:}
  !A \strictif !B,  !A \strictif !C & \Entails !A \strictif (!B \land !C)
  \intertext{কঠোর শর্তবচনের পূর্বপক্ষকে শক্তিশালী করা যায়:}
  !A \strictif !B & \Entails (!A \land !C) \strictif !B
  \intertext{কঠোর শর্তবচন পরিযায়ীও:}
  !A \strictif !B, !B \strictif !C & \Entails !A \strictif !C
  \intertext{শেষত, কঠোর শর্তবচন তার প্রতিবিপরীত রূপের সমতুল্য:}
  !A \strictif !B & \Entails \lnot !B \strictif \lnot !A\\
  \lnot !B \strictif \lnot !A & \Entails !A \strictif !B
\end{align}

\begin{prob}
  কঠোর শর্তবচনের ক্ষেত্রে যেসব অনুগমন সত্য নয়, সেগুলির
  জন্য \Log{S5}-এ প্রতিদৃষ্টান্ত দিন; অর্থাৎ নিচের দাবিগুলির
  জন্য:
  \begin{enumerate}
  \item $\lnot p \Entails/ \Box(p \lif q)$
  \item $q \Entails/ \Box(p \lif q)$
  \item $\lnot\Box(p \lif q) \Entails/ p \land \lnot q$
  \item $\Entails/ \Box(p \lif q) \lor \Box(q \lif p)$
  \end{enumerate}
\end{prob}

\begin{prob}
  বৈধ অনুগমনগুলি যে কঠোর শর্তবচনের জন্য সত্য, তা
  \Log{S5}-এ প্রমাণ দিয়ে দেখান:
  \begin{enumerate}
  \item $\Box(!A \lif !B) \Entails \lnot !A \lor !B$
  \item $!A \land \lnot !B \Entails \lnot\Box(!A \lif !B)$
  \item $!A, \Box(!A \lif !B) \Entails !B$
  \item $\Box(!A \lif !B), \Box(!A \lif !C) \Entails \Box(!A \lif (!B
    \land !C))$
  \item $\Box(!A \lif !B) \Entails \Box((!A \land !C) \lif !B)$
  \item $\Box(!A \lif !B), \Box(!B \lif !C) \Entails \Box(!A
    \lif !C)$
  \item $\Box(!A \lif !B) \Entails \Box(\lnot !B \lif \lnot !A)$
  \item $\Box(\lnot !B \lif \lnot !A) \Entails \Box(!A \lif !B)$
  \end{enumerate}
\end{prob}

তবে S5 ব্যবস্থায় কঠোর শর্তবচনেরও নিজস্ব ``আপাত-বৈপরীত্য'' আছে।
মিথ্যা পূর্বপক্ষ বা সত্য অনুপক্ষযুক্ত বস্তুগত শর্তবচন
যেমন সত্য হয়, তেমনই \emph{আবশ্যিকভাবে} মিথ্যা পূর্বপক্ষ
অথবা আবশ্যিকভাবে সত্য অনুপক্ষযুক্ত কঠোর শর্তবচনও সত্য।
এ ছাড়া, সত্য কঠোর শর্তবচন আবশ্যিকভাবে সত্য এবং মিথ্যা
কঠোর শর্তবচন আবশ্যিকভাবে মিথ্যা। অন্য কথায়:
\begin{align}
  \Box \lnot !A & \Entails !A \strictif !B\\
  \Box !B & \Entails !A \strictif !B\\
  !A \strictif !B& \Entails \Box(!A \strictif !B)\\
  \lnot(!A \strictif !B) & \Entails \Box\lnot(!A \strictif !B)
\end{align}
$\strictif$-কে ``নিহিত করে'' অর্থে নিলে এগুলি সমস্যা নয়।
যৌক্তিক অনুগমন তো গাণিতিক সত্য; তাই তা আপতিক হতে পারে না।
কিন্তু ইংরেজি ``if \dots then \dots''-এর অর্থ প্রকাশের
সংযোজক হিসেবে $\strictif$ ব্যবহার করতে চাইলে, বিশেষত
শেষের দুটি তথ্য সমস্যা তোলে। নিশ্চয়ই এমন শর্তবচন আছে
যেগুলি আপতিকভাবে সত্য বা আপতিকভাবে মিথ্যা; বস্তুত,
সাধারণত এগুলি আবশ্যিক সত্য কিংবা অসম্ভব—কোনোটিই নয়।

\begin{prob}
  \Log{S5}-এ নিচের দাবিগুলির প্রমাণ দিন:
  \begin{enumerate}
    \item $\Box \lnot !A \Entails !A \strictif !B$
    \item $!A \strictif !B \Entails \Box(!A \strictif !B)$
    \item $\lnot(!A \strictif !B)  \Entails \Box\lnot(!A \strictif !B)$
  \end{enumerate}
  প্রমাণে $\strictif$-এর সংজ্ঞা ব্যবহার করুন।
\end{prob}

\end{document}
